\subsection{Characters and selection of a vacuum}
\label{ex-phase-selection}

The preceding section constructs a commutative algebra on the whole ground
space. Its characters provide a useful way of separating the vacuum state
into components. The argument below also explains why a one-dimensional
ground space is first obtained almost everywhere, rather than assumed
for the unselected table. After relativistic reconstruction, a uniform
low-energy estimate extends vacuum uniqueness to every character.

\begin{proposition}[The representing probability]\label{ex-character-measure}
Suppose the algebra \(\mathcal A\) of
Proposition~\ref{ex-ground-commutative} is separable in operator norm.
Its space \(\Xi\) of unital complex star characters, with the topology
of pointwise convergence on \(\mathcal A\), is compact and metrizable.
There is a probability measure \(\mu\) on \(\Xi\), positive on every
nonempty open set, such that
\begin{equation}\label{ex-spectral-state}
 \int_\Xi\chi(A)\,d\mu(\chi)=\langle\Omega,A\Omega\rangle
 \quad(A\in\mathcal A).
\end{equation}
\end{proposition}
\begin{proof}
The algebra is a nonzero unital commutative \(C^*\)-algebra: it is a
norm-closed star subalgebra of \(B(\mathcal H_0)\), and
\(\|\Omega\|=1\). The commutative Gelfand theorem identifies it
isometrically, as a star algebra, with \(C(\Xi)\), by
\(A\mapsto\widehat A\), where \(\widehat A(\chi)=\chi(A)\)
\cite[Theorem 2.13]{Williams}. Its character space is compact since the
algebra is unital.

Choose a countable norm-dense sequence \((A_j)\) in \(\mathcal A\).
The continuous map \(\chi\mapsto(\chi(A_j))_j\) into
\(\mathbb C^{\mathbb N}\) is injective: equality on the dense sequence
implies equality everywhere by continuity of the characters. A continuous
injection of a compact space into a Hausdorff space is a homeomorphism onto
its image (images of closed sets are compact and hence closed).
The countable product is metrizable, for example by
\(\sum_j2^{-j-1}\min(1,|z_j-z'_j|)\), so \(\Xi\) is metrizable.

Transfer the functional \(A\mapsto\langle\Omega,A\Omega\rangle\)
to \(C(\Xi)\). It is positive: a nonnegative continuous function is
\(\overline g g\) for its continuous nonnegative square root \(g\),
and its inverse Gelfand transform has the form \(B^*B\). The functional
takes the value \(\|B\Omega\|^2\ge0\) there. The Riesz representation
theorem supplies a positive Borel measure representing it
\cite[Chapter 6, Theorem 32]{Sawyer}. Compactness gives
\(C_c(\Xi)=C(\Xi)\), so that theorem applies directly. At the constant
function one it gives \(\mu(\Xi)=\|\Omega\|^2=1\).

To prove full support, suppose a nonempty open set \(O\) has zero measure.
Choose \(x\in O\) and \(r>0\) with \(B(x,r)\subset O\). The function
\(f(y)=\max(0,r-\operatorname{dist}(x,y))\) is continuous, nonzero at
\(x\), and zero outside \(O\). Let \(A\) be its inverse Gelfand
transform. Then \eqref{ex-spectral-state} applied to \(A^*A\) gives
\(\|A\Omega\|^2=\int|f|^2\,d\mu=0\). Separation of \(\Omega\)
forces \(A=0\), contradicting \(f(x)=r\). Thus every such open set has
positive measure.
\end{proof}

For \(\chi\in\Xi\), define the whole-word values
\begin{equation}\label{ex-phase-word}
 S_\chi(w)=\chi(M(w))
 \quad\text{for distinct-point words }w.
\end{equation}
It is essential to apply the character to the operator for the entire
word. The compressed operator \(M(w_1w_2)\) need not equal
\(M(w_1)M(w_2)\) at finite separation.
Equations~\eqref{ex-spectral-state} and \eqref{ex-compression-values} give
the barycenter identity
\begin{equation}\label{ex-phase-barycenter}
 S(w)=\int_\Xi S_\chi(w)\,d\mu(\chi).
\end{equation}
For a fixed word the integrand is continuous, and hence integrable.

Some properties of \(S_\chi\) already follow at this level. Normalization
and translation invariance follow from \(M(\varnothing)=I\) and
\eqref{ex-compression-values}. It is also reflection positive. For
positive words \(v_i\) and coefficients \(c_i\), the Gram identity in
the preceding section gives
\[
 \sum_{i,j}\overline{c_i}c_jM((\theta v_i)v_j)
 =\left(\sum_i c_iX(v_i)\right)^*
   \left(\sum_j c_jX(v_j)\right)\ge0
 \quad\text{in }B(\mathcal H_0).
\]
The operator on the left lies in \(\mathcal A\). Positivity is the
same in a unital \(C^*\)-subalgebra as in the ambient operator algebra
(spectral permanence and continuous functional calculus give its nonnegative square root in
the subalgebra; see \cite[Theorems 2.17 and 2.19]{Williams}). Characters take positive elements to nonnegative
numbers. Applying \(\chi\) proves the desired Gram positivity.
Regularity and the remaining covariance and field identities still need
their own operator estimates; the barycenter formula alone does not
establish them.

\begin{lemma}[A common countable collection of words]\label{ex-phase-horizon}
Suppose the raw label spaces are finite-dimensional complex vector spaces
indexed by the nonnegative integers. Suppose that each \(S_\chi\) is
multilinear in its labels and continuous on the distinct-point
configuration space, and has the Hilbert realization of
Proposition~\ref{ex-gns}. There is one countable collection of positive
words whose vectors span a dense subspace in every \(\mathcal H_\chi\).
\end{lemma}
\begin{proof}
Choose a finite basis in every raw label space. Take all positive words
with rational coordinates and with basis-vector labels, for every arity
and every finite list of label grades; include the empty word. There are
countably many choices. Every word is a finite linear combination of
words with basis labels, by expanding each label in its basis. This identity holds for the Hilbert-space vectors: pair the difference with any word vector and use multilinearity of the correlation in the positive-word labels. All such inner products vanish, so density forces the difference to be zero.
For a fixed basis-labelled positive word, approximate its positions by
rational positions, remaining within the open set of positive
distinct-point configurations. Its vectors converge in every
\(\mathcal H_\chi\): the squared norm of their difference is the sum
of two diagonal Gram values minus twice the real part of the cross Gram
value, and continuity makes this tend to zero. The reflected and positive
halves remain strictly separated throughout this limit. The full word
span is dense by construction; thus the selected countable span is dense
as well. The positions and bases were chosen without reference to
\(\chi\), so the collection is common to every phase.
\end{proof}

\begin{theorem}[A one-dimensional ground space in almost every phase]
\label{ex-pure-ground}
Assume every \(S_\chi\) has the Euclidean Hilbert-space dynamics of
Theorem~\ref{ex-euclidean-semigroup}, with unit vacuum \(\Omega_\chi\).
Assume there is a common countable collection of positive words whose
vectors span densely in every phase Hilbert space. Then, for
\(\mu\)-almost every \(\chi\),
\[
 \mathcal H_{\chi,0}=\mathbb C\Omega_\chi.
\]
In particular there exists a phase with this property.
\end{theorem}
\begin{proof}
Fix a positive word \(v\). For integers \(n\ge0\), set
\[
 f_n(\chi)=\|E_{\chi,n}[v]_\chi\|^2
 =S_\chi((\theta\tau_nv)(\tau_nv)).
\]
Each is a continuous function of \(\chi\), by
\eqref{ex-phase-word}. Contractivity gives \(0\le f_n\le f_0\);
the latter is integrable on the compact probability space. The
large-time projection theorem gives
\(f_n(\chi)\to\|P_{\chi,0}[v]_\chi\|^2\).
Dominated convergence and the barycenter identity imply
\begin{equation}\label{ex-phase-ground-average}
 \int_\Xi\|P_{\chi,0}[v]_\chi\|^2\,d\mu
 =\lim_{n\to\infty}\|E_n[v]\|^2=\|P_0[v]\|^2.
\end{equation}
All limit functions here are measurable as pointwise limits of the
continuous \(f_n\).

Apply \eqref{ex-spectral-state} to \(M(v)^*M(v)\). Since characters
preserve products and adjoints and \(M(v)\Omega=P_0[v]\),
\[
 \int_\Xi|S_\chi(v)|^2\,d\mu=\|P_0[v]\|^2.
\]
Moreover \(\langle\Omega_\chi,P_{\chi,0}[v]_\chi\rangle=S_\chi(v)\),
because the projection fixes the vacuum. Orthogonal decomposition along
the unit vector \(\Omega_\chi\) gives
\[
 \|P_{\chi,0}[v]_\chi-S_\chi(v)\Omega_\chi\|^2
 =\|P_{\chi,0}[v]_\chi\|^2-|S_\chi(v)|^2.
\]
The nonnegative right side is integrable and, by the two preceding
identities, has integral zero. It is therefore zero almost everywhere.

Apply this conclusion to the common countable collection and intersect
its full-measure sets. For every \(\chi\) in the intersection, the
projection of its dense word span lies in \(\mathbb C\Omega_\chi\).
This one-dimensional subspace is closed. Continuity of the projection
therefore implies \(\operatorname{ran}P_{\chi,0}\subset
\mathbb C\Omega_\chi\). The reverse inclusion holds because the vacuum
is fixed. Finally a set of measure one in a probability space is nonempty.
\end{proof}

\begin{remark}[Only the ground space is used in the selection]
\label{rem:ex-selection-ground-only}
The character \(\chi_*\) is chosen from the full-measure set of
Theorem~\ref{ex-pure-ground}. No condition on the short-distance
behaviour of its correlations enters the choice. The trace asymptotic of
Theorem~\ref{thm:trace-asymptotic} is proved afterwards at the selected
character. The comparison bounds in Step~2 of
Subsection~\ref{subsec:trace-rounded-comparison} hold uniformly over all
characters, and the identification in its Step~3 uses the null relations at
the selected character. One could add a second condition to the selection:
for every \(e\in\mathbb R\), the product \(r^8\) times the two-point function
of the shifted trace \(\Theta+e\) tends to zero along a fixed subsequence of
radii. This would change nothing. That two-point function equals
\(T(r)+e^2\), because the one-point function of \(\Theta\) vanishes. The
conclusion \eqref{ex-trace-limit} gives \(r^8T(r)\to0\), so the condition
holds at \(\chi_*\) for every \(e\) and along every sequence \(r\downarrow0\).
\end{remark}

This selection argument requires the phase Hilbert spaces and their
contraction semigroups before it can be applied. It does not construct
them from a probability decomposition alone. Retention of the raw-label
null spaces, field derivatives, stress tensor and relativistic
reconstruction remain additional obligations for the chosen phase.

\subsubsection{Vacuum uniqueness at every character}

Almost-everywhere vacuum uniqueness does not by itself imply uniqueness
at every character: a positive spectral value could approach zero at an
exceptional character. The following quantitative argument excludes this
possibility. It uses the relativistic reconstruction of each character
\emph{without} its vacuum-uniqueness clause. That construction uses the
common phase caps and slab estimates of Subsection~\ref{ex-phase-regularity},
the continuation and boundary arguments, and the word construction.
Its forward spectrum follows from the first assertion of
Theorem~\ref{thm:field-spectrum}, which does not assume a unique vacuum.
Thus the use of that assertion below is not circular.

For a character $\chi$, write $U_\chi(t,x)$ for the reconstructed
translations, with $t\in\mathbb R$ and $x\in\mathbb R^3$, and let
$P_\chi$ be their momentum measure. Our convention is
$U_\chi(t,x)=\int e^{i(Et+p\cdot x)}\,dP_\chi(E,p)$.
Write $H_\chi$ for the nonnegative energy generator and
$Q_\chi=P_\chi(\{E=0\})$. Its range is the translation-invariant
space, since the forward cone intersects $\{E=0\}$ only at the origin.
The vacuum belongs to this range.

A compactly smeared word sum is a finite linear combination $A_\chi$
of products of fields with smooth compactly supported spacetime tests,
acting on the common invariant domain. The labels, coefficients and tests
are fixed independently of $\chi$. Its formal adjoint $A_\chi^*$ is
obtained by reversing each word and conjugating labels, tests and
coefficients. Both $A_\chi\Omega_\chi$ and
$A_\chi^*\Omega_\chi$ are well defined. All norms and scalar products
below are vacuum word correlations, so their dependence on $\chi$ is
controlled by the common boundary construction.

\begin{lemma}[A uniform low-energy estimate]
\label{lem:phase-uniform-low-energy}
For each compactly smeared word sum there is $D_A<\infty$, independent
of the character, such that its positive spectral measure
\[
 \mu_\chi(B)=\langle A_\chi\Omega_\chi,
                  P_\chi(B)A_\chi\Omega_\chi\rangle
\]
satisfies
\begin{equation}\label{eq:phase-low-energy}
 \mu_\chi\{0<E\le u\}\le D_Au^2\qquad(u>0).
\end{equation}
Consequently, for every $s>0$,
\begin{equation}\label{eq:phase-heat-remainder}
 0\le\langle A_\chi\Omega_\chi,
       (e^{-sH_\chi}-Q_\chi)A_\chi\Omega_\chi\rangle
       \le 2D_As^{-2}.
\end{equation}
\end{lemma}
\begin{proof}
Let $\nu_\chi$ be the analogous spectral measure for
$A_\chi^*\Omega_\chi$. The boundary continuity in Proposition~\ref{prop:assembly-boundary-barycenter} makes both total masses
continuous in $\chi$. Compactness of the character space gives one
$M<\infty$ bounding their sum. Both measures are supported in the
forward cone.

Choose $T,S<\infty$ so that every test support lies in
$\{|t|\le T,\ |x|\le S\}$, and take $R>1+2T+2S$.
For $|t|\le1$ and $|x|>R$, every point in a test support translated
by $(t,x)$ is spacelike separated from every point in any untranslated
test support: the spatial distance exceeds $|x|-2S>1+2T$, whereas
the time difference is at most $1+2T$. Locality allows all factors
in the translated word sum to commute past all factors of the
untranslated adjoint word sum on the invariant domain. Vacuum
translation invariance therefore gives
\[
 \delta_\chi(t,x):=
 \int e^{i(Et+p\cdot x)}\,d\mu_\chi(E,p)
 -\int e^{-i(Et+p\cdot x)}\,d\nu_\chi(E,p)=0.
\]
Indeed these two terms are respectively
$\langle\Omega_\chi,A_\chi^*A_\chi(t,x)\Omega_\chi\rangle$
and
$\langle\Omega_\chi,A_\chi(t,x)A_\chi^*\Omega_\chi\rangle$.
The same argument covers finite sums and the empty word.

Put
\[
 h(E)=\int_0^1\sin(tE)\,dt
 =\begin{cases}(1-\cos E)/E,&E>0,\\0,&E=0,\end{cases}
 \qquad
 g_\chi(x)=\int_0^1
       \frac{\delta_\chi(t,x)-\delta_\chi(-t,x)}{2i}\,dt.
\]
For $E\ge0$, $0\le h(E)\le1$; for $0\le E\le1$,
$\sin(tE)\ge2tE/\pi$ and hence $h(E)\ge E/\pi$.
Fubini is justified by the finite masses. It expresses $g_\chi$ as
the Fourier transform, with positive exponent, of the positive measure
\[
 \rho_\chi=(p)_*(h(E)\mu_\chi)
                    +(-p)_*(h(E)\nu_\chi)
             \quad\hbox{on }\mathbb R^3.
\]
The plus sign in the second term is important: time reversal changes
the sign of its sine as well as the sign preceding $\nu_\chi$.
We have $|g_\chi(x)|\le M$ and $g_\chi(x)=0$ for $|x|>R$.

For $0<u\le1$, choose a smooth $B_u:\mathbb R^3\to[0,1]$
equal to one on the ball of radius $u$ and supported in the ball of
radius $2u$. Define
$\widehat B_u(x)=\int e^{-ip\cdot x}B_u(p)\,dp$.
Then $|\widehat B_u(x)|\le\operatorname{vol}(B_{2u})$.
Fourier inversion and Fubini give
\[
 (2\pi)^3\int B_u\,d\rho_\chi
 =\int g_\chi(x)\widehat B_u(x)\,dx
 \le M\operatorname{vol}(B_R)\operatorname{vol}(B_2)u^3.
\]
For Fubini one may use the rapid decay of $\widehat B_u$ and the
finite mass of $\rho_\chi$. The left side is real and nonnegative;
the inequality follows by taking the absolute value of the right side.
On $u/2<E\le u$, cone support gives $|p|\le u$, so $B_u(p)=1$
and $h(E)\ge u/(2\pi)$. Positivity of both terms in $\rho_\chi$
therefore yields
\[
 \mu_\chi\{u/2<E\le u\}\le C_Au^2,
 \qquad
 C_A=\frac{M\operatorname{vol}(B_R)\operatorname{vol}(B_2)}{(2\pi)^2}.
\]
The disjoint dyadic rings with upper endpoints $2^{-j}u$ cover
$\{0<E\le u\}$. Summing their bounds gives at most
$(4/3)C_Au^2\le2C_Au^2$. For $u>1$ use
$\mu_\chi(\mathbb R^4)\le M\le Mu^2$.
Thus \eqref{eq:phase-low-energy} holds with $D_A=\max(M,2C_A)$.

Finally the spectral theorem identifies the middle expression of
\eqref{eq:phase-heat-remainder} with
$\int_{E>0}e^{-sE}\,d\mu_\chi$. Since
$e^{-sE}=\int_E^\infty se^{-st}\,dt$, Tonelli and
\eqref{eq:phase-low-energy} bound this by
\[
 s\int_0^\infty e^{-st}\mu_\chi\{0<E\le t\}\,dt
 \le D_As\int_0^\infty t^2e^{-st}\,dt=2D_As^{-2}.
\]
All constants depend on the fixed word sum, not on the character.
\end{proof}

\begin{theorem}[The vacuum line in every reconstructed character]
\label{thm:ex-every-character-vacuum}
For every character $\chi$, the translation-invariant vectors in its
reconstructed Hilbert space are exactly $\mathbb C\Omega_\chi$.
\end{theorem}
\begin{proof}
For almost every character, Theorem~\ref{ex-pure-ground} and the
heat-vector isometry imply this assertion by the last part of
Theorem~\ref{thm:field-spectrum}. Fix a compactly smeared word sum and
$s>0$. The scalar function
\[
 \chi\longmapsto
 \langle A_\chi\Omega_\chi,e^{-sH_\chi}A_\chi\Omega_\chi\rangle
       -|\langle\Omega_\chi,A_\chi\Omega_\chi\rangle|^2
\]
is continuous by Proposition~\ref{prop:assembly-boundary-barycenter}:
the heat pairing is the same vacuum word distribution at a fixed
nonnegative imaginary gap $s$ between its two halves. Indeed the
mixed Gram identity \eqref{eq:point-gram-continuation}, integrated
against the tests, gives this identity at strictly positive internal
heights. Send those heights to zero using the independent-height
vector limits of Subsection~\ref{sec:heat-vectors}. The insertion of
$e^{-sH_\chi}$ is a bounded contraction and commutes with the
translation intertwiner, so the vector pairings converge. The common
weakly ordered boundary limit gives the stated fixed-gap distribution. On the
full-measure set just obtained, $Q_\chi$ is the vacuum projection, so
\eqref{eq:phase-heat-remainder} places this function in the closed
interval $[0,2D_As^{-2}]$. The bound holds everywhere. Otherwise
continuity would give a nonempty open set on which it fails, and full
support of $\mu$ would give this set positive measure, a contradiction.

Fix now any character and let $s\to\infty$. Spectral dominated
convergence gives $e^{-sH_\chi}\to Q_\chi$ strongly. The preceding
bound implies
\[
 \|Q_\chi A_\chi\Omega_\chi\|^2
       =|\langle\Omega_\chi,A_\chi\Omega_\chi\rangle|^2.
\]
Since $Q_\chi\Omega_\chi=\Omega_\chi$, orthogonal decomposition along
the unit vector $\Omega_\chi$ shows
\[
 Q_\chi A_\chi\Omega_\chi
=\langle\Omega_\chi,A_\chi\Omega_\chi\rangle\Omega_\chi.
\]

Such compactly smeared word vectors span densely. To verify the compact
support restriction, choose a smooth cutoff $\zeta=1$ near zero and
replace each Schwartz test $f$ by $\zeta(\,\cdot/n)f$. These converge
to $f$ in every Schwartz seminorm, by Leibniz' rule and rapid decay on
the escaping support of the cutoff derivatives. Joint continuity of
word correlations implies convergence of the resulting word vectors:
the squared norm difference is a finite sum of convergent Gram values.
Cyclicity then gives density. Since $Q_\chi$ is bounded and the vacuum
line is closed, its range on the whole Hilbert space lies in that line.
The reverse inclusion follows from vacuum invariance.
\end{proof}
